\documentclass[11pt]{amsart}
\usepackage{amsmath,amssymb,amsthm}
\usepackage[margin=1.1in]{geometry}
\usepackage{hyperref}
\usepackage{booktabs}

\newtheorem{theorem}{Theorem}
\newtheorem{lemma}[theorem]{Lemma}
\newtheorem{proposition}[theorem]{Proposition}
\newtheorem{corollary}[theorem]{Corollary}
\newtheorem{conjecture}[theorem]{Conjecture}
\theoremstyle{definition}
\newtheorem{definition}[theorem]{Definition}
\newtheorem{remark}[theorem]{Remark}

\newcommand{\Ev}{\mathrm{Ev}}
\newcommand{\Od}{\mathrm{Od}}
\newcommand{\tr}{\to}

\title[Five even vertices with one universal vertex]{Five even vertices with one universal vertex:\\ Corollary~106 without the unrefereed input}
\author{Justin Williams}
\date{November 2026 (working draft)}

\begin{document}
\begin{abstract}
Corollary~106 of the marked-reduction manuscript~\cite{MR} asserts a floor bound for connected
graphs in which at most one vertex of the subgraph induced by the even-degree vertices has degree at
least four. Its one-exception case at odd order rested on an unrefereed preprint~\cite{Shaik}. The
natural replacement suggested in the project notes, a version of Botler and Sambinelli's
theorem~\cite[Theorem~13]{BS} with one prescribed end, is false: we give counterexamples from
order five on, an infinite family with a prescribed vertex of degree two between two odd vertices,
and an order-nine graph on which the doubling gadget of the notes fails although the floor holds.

We then show that the corollary is only ever used with at most five unmarked even vertices, and we
prove that case directly. \emph{Let $W$ have odd order $N$ and exactly five even vertices, which
induce a graph with exactly one vertex adjacent to the other four. Then $p(W)\le(N-1)/2$, provided
either the even subgraph is not a bowtie, or $W$ is connected and has a vertex of degree one.} In
fact the even vertices then admit a realisable pair. The proof uses only Lov\'asz's theorem and
Fan's addible-edge lemmas. Six of the seven possible even subgraphs are settled by three short
lemmas that need no leaf. The bowtie is settled by eleven explicit transfer certificates of at most
six moves each. Consequently Corollary~106 holds without~\cite{Shaik} in every application in the
manuscript, and so does Theorem~144, the reduction of the degree-six deletion test to five even
vertices. All certificates are verified by an independent script and replayed in Lean~4 on an
arbitrary host graph.
\end{abstract}
\maketitle

\section{Introduction}

All graphs are finite and simple. A \emph{partition} of a graph is a partition of its edge set into
paths of positive length, and $p(G)$ is the least size of one. For a partition $D$ and a vertex $v$
let $b_D(v)$ be the number of ends of paths of $D$ at $v$. Then $b_D(v)\equiv d(v)\pmod 2$ and
$\sum_v b_D(v)=2|D|$. We write $\Ev=\Ev(G)$ for the set of vertices of positive even degree,
$\Od$ for the odd vertices, and $G[\Ev]$ for the \emph{even subgraph}. A set $S\subseteq\Ev$ is
\emph{realisable} if some partition has one end at every odd vertex, two at every vertex of $S$ and
none at $\Ev\setminus S$; such a partition has $|\Od|/2+|S|$ paths~\cite{ES}.

\subsection*{The gap}
Corollary~106 of~\cite{MR} reads: \emph{let $G$ be connected and $S$ a nonempty set of marks; if at
most one vertex of the subgraph induced by the unmarked even vertices has degree at least four in
that subgraph, then the marked bound $q(G,S)\le\lfloor(|V(G)|-|S|)/2\rfloor$ holds.} Its proof
reduces everything to a \emph{leaf floor bound}: if $L$ is connected, has a leaf at an odd support,
and $L[\Ev(L)]$ has at most one vertex of degree at least four, then $p(L)\le\lfloor|L|/2\rfloor$.
With no exceptional vertex this is Botler--Sambinelli~\cite[Theorem~13]{BS} (the leaf excludes their
exceptional SET graphs). With one exceptional vertex and $|L|$ even it is Fan--Hou--Zhou's
Theorem~5~\cite{FHZ}, since ceiling and floor agree. The remaining case, one exceptional vertex and
$|L|$ odd, was taken from Shaik's preprint~\cite{Shaik} by doubling $L$ along a bridge, which creates
two exceptional vertices (audit~A7, \S4). This is the last unrefereed input of Theorem~144 of~\cite{MR}.

\subsection*{The prescribed-end route fails}
The project notes (connections~\S L) observe that every gadget that turns the doubled host into a
graph covered by~\cite{BS} loses exactly one unit of budget, and that the missing ingredient would
be Botler--Sambinelli's theorem with one prescribed end: \emph{if $N$ is odd and the even subgraph
has maximum degree at most three, then for every even $y$ there is a partition into $(N-1)/2$ paths
with an end at $y$.} Section~\ref{sec:false} shows that this is false in a strong sense:
\begin{itemize}
\item it fails at every odd order $N\ge5$ and for every odd number $e\ge3$ of even vertices, for an
infinite family in which $y$ has degree two, both neighbours of $y$ are odd, and the even subgraph
is a matching together with the isolated vertex $y$ (Proposition~\ref{prop:family}). This is
exactly the configuration of the joining gadget $x-y-x'$;
\item exhaustively at $N=7$ there are $6$ failures with five even vertices
(Table~\ref{tab:census});
\item the gadget itself fails on a graph $L$ of order nine that satisfies the floor bound and lies
in the case needed by the manuscript (Proposition~\ref{prop:gadget}).
\end{itemize}
The obstruction in every example is the same elementary forcing lemma (Lemma~\ref{lem:forcing}):
a cycle whose vertices other than one have degree two must contain an end. Prescribing $y$ adds
one forced end to a configuration that may already be tight.

\subsection*{What Corollary~106 needs, and the main result}
Every use of Corollary~106 in~\cite{MR} has at most five unmarked even vertices
(Section~\ref{sec:needs}). For five unmarked even vertices a single exceptional vertex has degree
four, so it is \emph{universal} in the five-vertex even subgraph. This is the case we prove.

\begin{theorem}\label{thm:main}
Let $W$ be a graph of odd order $N$ without isolated vertices, with exactly five even vertices
$T=\Ev(W)$, such that $K=W[T]$ has exactly one vertex adjacent to all other vertices of $T$.
\begin{enumerate}
\item[(a)] If $K$ is not a bowtie (two triangles sharing one vertex), some $2$-subset of $T$ is
realisable. In particular $p(W)\le(N-1)/2$.
\item[(b)] If $K$ is a bowtie, the same holds provided $W$ has adjacent odd vertices $u,v$ such that
$u$ has a neighbour in $T$ and $v$ has none. Such $u,v$ exist when $W$ is connected and has a vertex
of degree one whose neighbour is odd.
\end{enumerate}
\end{theorem}

\begin{corollary}\label{cor:bp5}
Let $L$ be a connected graph of odd order with a vertex of degree one and exactly five even
vertices, whose even subgraph has at most one vertex of degree four. Then
$p(L)\le\lfloor|L|/2\rfloor$.
\end{corollary}

\begin{proof}
With no vertex of degree four the even subgraph is subcubic, and~\cite[Theorem~13]{BS} applies: a
SET graph has exactly three even vertices. With one such vertex and the leaf at an odd support,
Theorem~\ref{thm:main} applies. If the leaf $\ell$ is at an even support $r$, then $L-\ell$ is
connected of even order with exactly four even vertices, and~\cite[Corollary~22]{ES} (Lov\'asz and
Fan only; Lean \texttt{rooted\_four}) gives a partition of $L-\ell$ into $|L-\ell|/2$ paths with an
end at the now odd vertex $r$, which absorbs the leaf edge.
\end{proof}

Corollary~\ref{cor:bp5} is exactly the leaf floor bound in the one-exception odd-order case with
five even vertices. Section~\ref{sec:cor106} states the corrected Corollary~106 and checks its three
uses in Theorem~144, the consequence for Proposition~105, and the remark after Theorem~114. The
one-exception odd-order case with seven or more unmarked even vertices remains open without the
preprint, but nothing in~\cite{MR} uses it.

\subsection*{Method}
The proof is a set of \emph{Lov\'asz--Fan certificates}: hold a few edges among the even vertices
(and at most two exterior vertices), so that exactly one vertex of the remaining graph has even
degree; take the exact Lov\'asz partition; and restore the held edges one at a time by Fan's single
or grouped transfer. This is the certificate format of~\cite{ES} and of the audit's seven-site search
(A6 \S3.2), which had certified all seven cores by exhaustive computation. Here six cores are
handled by three general lemmas (Section~\ref{sec:local}), and the bowtie by a table of eleven
certificates (Section~\ref{sec:bowtie}). The scripts in \texttt{scripts/} regenerate and check every
certificate, and \texttt{lean/GallaiLean/FiveEven.lean} with \texttt{FiveEvenCerts.lean} replays
all of them on an arbitrary host (Section~\ref{sec:lean}).

\section{Tools}\label{sec:tools}

\begin{lemma}[Lov\'asz~\cite{L68}, exact form]\label{lem:lovasz}
If at most one vertex of $K$ has positive even degree, then $K$ has a partition with exactly one end
at every odd vertex and none elsewhere.
\end{lemma}

Apply Lov\'asz's theorem to each component. A component without even vertices has all its vertices
odd, so its $|C|/2$ paths have one end at each vertex. The component containing the even vertex has
odd order, and its $(|C|-1)/2$ paths put their $|C|-1$ ends on its $|C|-1$ odd vertices. This form
is \texttt{lovasz\_ends} in Lean, proved from Fan's lemma~\cite{LF}.

For a partition $D$ of a graph $K$ and a vertex $r$, a \emph{blank neighbour} of $r$ is a neighbour
$w$ of $r$ in $K$ with $b_D(w)=0$.

\begin{lemma}[Fan~\cite{Fan}, Lemmas 3.3 and 3.4]\label{lem:fan}
Let $D$ be a partition of $K$, let $r$ be a vertex and let $T'$ be a set of non-neighbours of $r$
with $r\notin T'$. If $\sum_{t\in T'}b_D(t)$ exceeds the number of blank neighbours of $r$, then for
some $t\in T'$ with $b_D(t)\ge1$ the graph $K+tr$ has a partition $D'$ with $|D'|=|D|$,
$b_{D'}(r)=b_D(r)+1$, $b_{D'}(t)=b_D(t)-1$, and $b_{D'}=b_D$ elsewhere.
\end{lemma}

Both forms are proved in Lean (\texttt{fan\_single}, \texttt{fan\_grouped}). We write
$t\tr r$ for the single transfer ($T'=\{t\}$) and $\{t,t'\}\tr r$ for the grouped transfer, after
which the argument continues with whichever donor the lemma returns. Every move adds one edge and
keeps the number of paths. A move is legal when the \emph{supply} $\sum_{t\in T'}b_D(t)$ is larger
than the number of \emph{blanks}, the blank neighbours of $r$.

\begin{definition}
A \emph{certificate} for a graph $W$ consists of a set $H\subseteq E(W)$ of \emph{held} edges such that
exactly one vertex of $W-H$ has even degree, and a sequence (a tree, for grouped moves) of legal moves
that restores all of $H$, starting from the partition of Lemma~\ref{lem:lovasz} for $W-H$.
\end{definition}

\begin{lemma}\label{lem:count}
If $W$ has odd order $N$, no isolated vertex, and a certificate, then $p(W)\le(N-1)/2$. If moreover
the final end vector is $1$ on $\Od(W)$, $2$ on two vertices of $T=\Ev(W)$ and $0$ on the others, the
final partition realises that pair.
\end{lemma}

\begin{proof}
By Lemma~\ref{lem:lovasz} the seed has $N-1$ ends: all vertices except the even one have one end.
Moves preserve the number of paths, hence the sum of ends.
\end{proof}

In all our certificates the held edges join vertices of a fixed \emph{local set} $M$ (the five even
vertices, sometimes with two odd vertices $u,v$). Every vertex outside $M$ keeps its degree, so it is
odd in $W-H$, has one end in the seed, and is never a donor or a receiver. Hence \emph{a vertex
outside $M$ is never blank}. The legality of each move then depends only on the adjacency inside $M$
and on the end counts on $M$. A certificate written on $M$ is therefore valid for every host $W$ with
the same adjacency inside $M$ and the same parities. Holding the edges may disconnect $W-H$; this is
allowed, since Lemma~\ref{lem:lovasz} holds componentwise.

\section{The prescribed-end lemma is false}\label{sec:false}

\begin{lemma}[Forcing]\label{lem:forcing}
Let $C$ be a cycle of $G$ through a vertex $w$, and suppose every vertex of $C$ other than $w$ has
degree two in $G$. Then every partition of $G$ has an end at some vertex of $C-w$.
\end{lemma}

\begin{proof}
At a vertex of degree two that is not an end, both edges lie on the same path. If no vertex of
$C-w$ were an end, all edges of $C$ would lie on one path. That path would then contain a cycle.
\end{proof}

\begin{proposition}\label{prop:family}
For integers $t,t'\ge0$ with $t+t'\ge1$, let $G_{t,t'}$ consist of a vertex $y$ adjacent to two
vertices $s,s'$, together with $t$ triangles $su_iv_i$ and $t'$ triangles $s'u_j'v_j'$ (pendant
triangles; if $t'=0$ then $s'$ is a leaf, and likewise for $t=0$). Then:
\begin{itemize}
\item $G_{t,t'}$ is connected of odd order $N=3+2(t+t')$;
\item its even subgraph is the perfect matching on the triangle vertices together with the
isolated vertex $y$;
\item $y$ has degree two, and both its neighbours are odd.
\end{itemize}
Nevertheless no partition of $G_{t,t'}$ into $(N-1)/2$ paths has an end at $y$.
\end{proposition}

\begin{proof}
The vertices $s,s'$ have degrees $1+2t$ and $1+2t'$. All other vertices have degree two. So the
only odd vertices are $s,s'$. By Lemma~\ref{lem:forcing} every partition has an end at one of $u_i,v_i$
for each triangle, and that vertex is even, so it has at least two ends. If $y$ is also an end, the
number of ends is at least $2+2(t+t'+1)=N+1$, which is more than $N-1$.
\end{proof}

The family covers every odd order $N\ge5$ and every odd number $e=1+2(t+t')\ge3$ of even vertices.
It is exactly the configuration of the joining gadget $x-y-x'$, with $s,s'$ in the roles of $x,x'$.
Each $G_{t,t'}$ satisfies the floor bound by~\cite[Theorem~13]{BS} (it is not a SET graph). So the
obstruction comes from the prescription alone: prescribing $y$ adds one forced end to a
configuration that is already tight.

At $N=7$ there are exactly six pairs $(G,y)$ with five even vertices and a subcubic even subgraph
for which the prescription fails (Table~\ref{tab:census}). They are:
\begin{itemize}
\item $G_{1,1}$ and $G_{2,0}$;
\item $y$ the attachment vertex of a pendant triangle, adjacent to a leaf and to an odd vertex that
carries a second pendant triangle;
\item $y$ of degree four in a pendant $K_4-e$, adjacent to an odd vertex with a pendant triangle;
\item $y$ between a leaf and a vertex of a $K_{1,1,3}$, or between a leaf and a vertex of a $K_5$,
whose other vertices are the four other even vertices.
\end{itemize}

\begin{table}[ht]
\centering
\begin{tabular}{lrrrr}
\toprule
& $e=1$ & $e=3$ & $e=5$ & $e=7$\\
\midrule
pairs $(G,y)$, $N=5$ & 7 & 30 & 5 & -- \\
\quad no floor partition ending at $y$ & 7 & 9 (5 SET) & 0 & -- \\
\quad lost by the odd-order Fan game & 7 & 14 & 0 & -- \\
pairs $(G,y)$, $N=7$ & 130 & 1236 & 845 & 7 \\
\quad no floor partition ending at $y$ & 130 & 101 (29 SET) & 6 & 0 \\
\quad lost by the odd-order Fan game & 130 & 230 & 64 & 0 \\
\bottomrule
\end{tabular}
\caption{Connected graphs of order $N$ whose even subgraph has maximum degree at most three, and
even vertices $y$ ($e$ is the number of even vertices). For $e=1$ the prescription is impossible by
counting. First row: exact (\texttt{scripts/exact\_prescribed.py}, CP-SAT). Last row: Lov\'asz--Fan
certificates with at most four held edges (\texttt{scripts/odd\_fan\_game.py}). Raising the bound to
eight held edges recovers only one of the 64 pairs with $e=5$.}\label{tab:census}
\end{table}

The last row records a second negative finding. At even order the Fan game wins every root of every
graph with subcubic even subgraph through order eight (experiments E13). At odd order with a
prescribed end, Lov\'asz--Fan certificates are far from complete even in that regime: most pairs
that the game loses do have the required partition.

\begin{proposition}\label{prop:gadget}
Let $L$ have vertices $x,a_1,b_1,a_2,b_2,s,w,z,\ell$ and edges $xa_i$, $xb_i$, $a_ib_i$ ($i=1,2$),
$xs$, $xw$, $s\ell$, $sw$, $wz$. Then $L$ is connected of order nine, has a leaf $\ell$ at the odd
support $s$, its even vertices are $x,a_1,b_1,a_2,b_2$ and induce a bowtie, and $p(L)=4$. Let $D$
be the joining gadget of connections \S L: two disjoint copies of $L-\ell$, the bridge $ss'$, and a
new vertex $y$ adjacent to $x$ and $x'$. Then $D$ has order $17$, its even subgraph consists of four
disjoint edges and the isolated vertex $y$, and no partition of $D$ into $8$ paths has an end at $y$.
\end{proposition}

\begin{proof}
For $L$, the paths $a_1b_1xb_2a_2$, $a_1xa_2$, $\ell sxwz$ and $sw$ form a partition into four paths,
and four is the parity bound for its four odd vertices $s,w,z,\ell$. In $D$ the vertices $x,x'$ have
odd degree $7$, so the even vertices are $y$ and the eight vertices $a_i,b_i,a_i',b_i'$, and $D$ has
eight odd vertices. Each of the four triangles $xa_ib_i$, $x'a_i'b_i'$ forces an end at one of its
two even vertices (Lemma~\ref{lem:forcing}). A partition with an end at $y$ therefore has at least
$8+2\cdot5=18$ ends, that is at least nine paths.
\end{proof}

So the gadget cannot complete the argument even in the case Corollary~106 needs. The same forcing
shows that the gadget fails whenever the exceptional vertex carries enough pendant triangles, while
$L$ itself is unaffected.

\section{Where Corollary 106 is used}\label{sec:needs}

The corollary is cited three times, all in the proof of Theorem~144 of~\cite{MR}, and its second
sentence is used once.
\begin{enumerate}
\item \emph{``By Corollary~106 with mark $r$, at least two vertices of $K$ have degree four.''} Here
$W$ has odd order, a leaf at the odd support $r$, and exactly five even vertices $T$, and the claim
is $p(W)\le\lfloor|W|/2\rfloor$ when $K=W[T]$ has at most one universal vertex. This is
Corollary~\ref{cor:bp5}.
\item \emph{The graph $H=B-rw$ with four even vertices,} and \emph{the matching-cut shores with at
most four unmarked even vertices.} With at most four unmarked even vertices the unmarked even
subgraph is subcubic, so only the zero-exception case is used:~\cite[Theorem~13]{BS} (audit A6, F18).
\item \emph{The second sentence of Corollary~106} (in the two-even-support failure of
Proposition~105, two unmarked even sites have degree four) is used only in the next paragraph, for
$H$ with exactly seven even sites and two even supports, so with five unmarked even sites.
\end{enumerate}
The remark after Theorem~114 (``Shaik's ordinary bound also applies if $K$ has at most one vertex of
degree greater than three\ldots\ at $d_J(z)=6$ a failure requires two universal vertices'') concerns
the five-vertex graph $K=J[N(z)\setminus\{r\}]$. Both deletion tests at $d_J(z)=6$ are floor bounds
for $J-z$ with one or two private leaves at $r$. These graphs have odd order, leaves at the odd
support $r$, and even set $N(z)\setminus\{r\}$. So the remark also follows from
Corollary~\ref{cor:bp5}. The step ``$r\in T$'' of Theorem~144, which cited~\cite{Shaik} directly, needs
only~\cite[Theorem~13]{BS} or~\cite[Corollary~22]{ES} (A7, U4).

\section{Six cores without a leaf}\label{sec:local}

Throughout this section and the next, $W$ is as in Theorem~\ref{thm:main}. Write $x$ for the
universal vertex of $K$ and $R=T\setminus\{x\}$ for the \emph{rim}. Since $x$ is the only vertex of
degree four in $K$, every rim vertex has at most two neighbours in $R$. Up to isomorphism $K[R]$ is
one of
\[
4K_1,\quad K_2+2K_1,\quad P_3+K_1,\quad K_3+K_1,\quad 2K_2,\quad P_4,\quad C_4 .
\]
For $K[R]=2K_2$ the graph $K$ is the bowtie. The lemmas below hold for every graph of odd order with
exactly five even vertices, whether or not one of them is universal. Each one is a certificate as in
Section~\ref{sec:tools}, with local set $T$.

\begin{lemma}[Two disjoint even edges]\label{lem:match}
Let $T=\{p,q,p',q',z\}$ with $pq,p'q'\in E$, and suppose $p\not\sim z$, $p'\not\sim z$ and
$p'\not\sim q$. Then $\{p,p'\}$ is realisable.
\end{lemma}

\begin{proof}
Hold $H=\{pq,p'q'\}$. In $W-H$ the vertices $p,q,p',q'$ are odd and $z$ is the only even vertex.
In the seed only $z$ is blank. Move $q\tr p$: the only possible blank neighbour of $p$ is $z$, and
$p\not\sim z$, so the supply $1$ exceeds $0$ blanks. Now $q$ is blank. Move $q'\tr p'$: the possible
blank neighbours $z,q$ are not adjacent to $p'$. The final end vector is $2$ at $p,p'$, $0$ at
$q,q',z$ and $1$ elsewhere.
\end{proof}

\begin{lemma}[Odd star]\label{lem:oddstar}
Let $a\in T$ have exactly three neighbours $w,p,q$ in $T$, and let $z$ be the fifth even vertex.
Suppose $p$ is not adjacent to all of $w,q,z$, and $q$ is not adjacent to all of $w,p,z$. Then $\{a,p\}$
or $\{a,q\}$ is realisable.
\end{lemma}

\begin{proof}
Hold $H=\{aw,ap,aq\}$. In $W-H$ the vertex $a$ has degree $d(a)-3$, so it is odd, as are $w,p,q$.
The only even vertex is $z$, and $a\not\sim z$.
\begin{enumerate}
\item $w\tr a$ with supply $1$ and $0$ blanks. Now $b(a)=2$ and $w$ is blank.
\item $\{p,q\}\tr a$ with supply $2$ and at most $1$ blank (namely $w$). Say $t$ is returned, and
let $t'$ be the other one. Now $b(a)=3$ and $t$ is blank.
\item $a\tr t'$ with supply $3$. The blank neighbours of $t'$ lie in $\{w,t,z\}$, and $t'$ misses
one of them, so there are at most $2$.
\end{enumerate}
Finally $b(a)=b(t')=2$, the vertices $w,t,z$ have no end, and every odd vertex has one.
\end{proof}

\begin{lemma}[Universal star]\label{lem:unistar}
Let $x\in T$ be adjacent to $a,b,c\in T$, and let $d$ be the fifth even vertex. Suppose $c$ has no
neighbour in $T$ other than $x$, $a$ is not adjacent to both $b$ and $d$, and $b$ is not adjacent to
both $a$ and $d$. Then $\{a,c\}$ or $\{b,c\}$ is realisable.
\end{lemma}

\begin{proof}
Hold $H=\{xa,xb,xc\}$. Then $x,a,b,c$ are odd in $W-H$ and $d$ is the only even vertex.
\begin{enumerate}
\item $\{a,b\}\tr x$ with supply $2$ and at most $1$ blank ($d$). Say $t$ is returned, and let $t'$
be the other one. Now $b(x)=2$ and $t$ is blank.
\item $x\tr t'$ with supply $2$. The edge $xt'$ is still held, so the blank neighbours of $t'$ lie in
$\{t,d\}$, and $t'$ misses one of them.
\item $x\tr c$ with supply $1$. The only $T$-neighbour of $c$ is $x$, and $xc$ is held, so $c$ has
no blank neighbour.
\end{enumerate}
Finally $b(t')=b(c)=2$ and $b(x)=b(t)=b(d)=0$.
\end{proof}

\begin{proposition}\label{prop:six}
If $K[R]\neq2K_2$, some $2$-subset of $T$ is realisable.
\end{proposition}

\begin{proof}
Label the rim $a,b,c,d$.
\begin{itemize}
\item $4K_1$: Lemma~\ref{lem:unistar} with the universal vertex $x$, where $a,b,c,d$ are pairwise
nonadjacent.
\item $K_2+2K_1$ with rim edge $ab$: Lemma~\ref{lem:match} with $(p,q,p',q',z)=(a,b,c,x,d)$.
\item $P_3+K_1$ with rim path $abc$, and $P_4$ with rim path $abcd$: Lemma~\ref{lem:match} with
$(p,q,p',q',z)=(b,c,a,x,d)$. Here $b\not\sim d$, $a\not\sim d$ and $a\not\sim c$.
\item $K_3+K_1$ with rim triangle $abc$: Lemma~\ref{lem:oddstar} at $a$, with $(w,p,q,z)=(x,b,c,d)$;
neither $b$ nor $c$ is adjacent to $d$.
\item $C_4$ with rim cycle $abcd$: Lemma~\ref{lem:oddstar} at $a$, with $(w,p,q,z)=(x,b,d,c)$;
$b\not\sim d$.\qedhere
\end{itemize}
\end{proof}

This proves Theorem~\ref{thm:main}(a). No hypothesis on $W$ outside $T$ is used, not even
connectivity. For the bowtie, an exhaustive search over all held sets inside $T$ finds no
certificate of this kind (\texttt{scripts/certgen.py}).

\section{The bowtie}\label{sec:bowtie}

Let $K$ be the bowtie with centre $x$ and triangles $xab$, $xcd$. Let $u,v$ be adjacent odd vertices
with $A=N(u)\cap T\neq\emptyset$ and $N(v)\cap T=\emptyset$. The local set is $M=T\cup\{u,v\}$, and
its edges are those of $K$, the edges from $u$ to $A$, and $uv$. As explained in
Section~\ref{sec:tools}, a certificate on $M$ is valid for every such host.

\begin{lemma}\label{lem:uv}
If $W$ is connected, $T\neq\emptyset$, and $W$ has a vertex $\ell$ of degree one whose neighbour is
odd, then such $u,v$ exist.
\end{lemma}

\begin{proof}
Let $\ell=p_0,p_1,\dots,p_m$ be a shortest path from $\ell$ to $T$. Then $p_0\notin T$, and $p_1$ is
odd, so $p_1\notin T$; hence $m\ge2$. Put $u=p_{m-1}$ and $v=p_{m-2}$. By minimality neither is in
$T$, so both are odd, and $v$ has no neighbour in $T$. Moreover $u\sim p_m\in T$.
\end{proof}

The symmetry group of the bowtie fixing $x$ has order eight, generated by $(a\,b)$, $(c\,d)$ and
$(a\,c)(b\,d)$. It splits the $31$ nonempty contact sets into eleven orbits. Table~\ref{tab:bowtie}
gives a certificate for one representative of each. A move is written $t\tr r$ with
(supply$\,>\,$blanks). After a grouped move, each branch is listed separately and is labelled by the
returned donor. Every certificate ends with $1$ at $u$ and $v$ and with a realisable pair $S$.

\begin{table}[ht]
\centering\footnotesize
\setlength{\tabcolsep}{4pt}
\begin{tabular}{llp{9.3cm}}
\toprule
$A$ & held, even seed vertex & moves \\
\midrule
$x$ & $ab,cd,xu$; $u$ & $b\tr a\,(1{>}0)$, $d\tr c\,(1{>}0)$, $x\tr u\,(1{>}0)$; $S=ac$\\
$a$ & $xb,cd,au$; $u$ & $b\tr x\,(1{>}0)$, $d\tr c\,(1{>}0)$, $a\tr u\,(1{>}0)$; $S=xc$\\
$ac$ & $xb,cd,au$; $u$ & $b\tr x\,(1{>}0)$, $a\tr u\,(1{>}0)$, $d\tr c\,(1{>}0)$; $S=xc$\\
$ab$ & $xa,cd,bu,uv$; $v$ & $d\tr c\,(1{>}0)$, $x\tr a\,(1{>}0)$, $b\tr u\,(1{>}0)$, $u\tr v\,(2{>}0)$; $S=ac$\\
$abc$ & $xa,bu,cu$; $d$ & $x\tr a\,(1{>}0)$, $c\tr u\,(1{>}0)$, $u\tr b\,(2{>}1)$; $S=ab$\\
$xac$ & $xb,au,cu$; $d$ & $u\tr a\,(1{>}0)$, $c\tr u\,(1{>}0)$, $x\tr b\,(1{>}0)$; $S=ab$\\
$xabc$ & $xa,bu,cu$; $d$ & $u\tr b\,(1{>}0)$, $c\tr u\,(1{>}0)$, $x\tr a\,(1{>}0)$; $S=ab$\\
$xa$ & $xb,xc,xu,au$; $d$ & $u\tr a\,(1{>}0)$, $\{b,c\}\tr x\,(2{>}1)$, then\\
& & \quad[$b$]: $x\tr c\,(2{>}1)$, $x\tr u\,(1{>}0)$; $S=ac$;
[$c$]: $x\tr b\,(2{>}0)$, $x\tr u\,(1{>}0)$; $S=ab$\\
$abcd$ & $xa,bu,cu,du,uv$; $v$ & $x\tr a\,(1{>}0)$, $b\tr u\,(1{>}0)$, $\{c,d\}\tr u\,(2{>}1)$, then\\
& & \quad[$c$]: $u\tr d\,(3{>}2)$, $u\tr v\,(2{>}0)$; $S=ad$;
[$d$]: $u\tr c\,(3{>}2)$, $u\tr v\,(2{>}0)$; $S=ac$\\
$xab$ & $cd,xu,au,bu,uv$; $v$ & $d\tr c\,(1{>}0)$, $x\tr u\,(1{>}0)$, $\{a,b\}\tr u\,(2{>}1)$, then\\
& & \quad[$a$]: $u\tr b\,(3{>}2)$, $u\tr v\,(2{>}0)$; $S=bc$;
[$b$]: $u\tr a\,(3{>}2)$, $u\tr v\,(2{>}0)$; $S=ac$\\
$xabcd$ & $xu,au,bu,cu,du,uv$; $v$ & see the text\\
\bottomrule
\end{tabular}
\caption{Certificates for the bowtie, one per orbit of contact sets $A=N(u)\cap T$. The seed has no
end at the listed even vertex and one end at every other vertex.}\label{tab:bowtie}
\end{table}

\emph{The case $A=T$.} Hold all six edges at $u$ inside $M$. Then $u$ remains odd, every vertex of
$T$ becomes odd, and $v$ is the only even vertex of $W-H$ (it is isolated if $\ell=v$ is a leaf). The
moves are:
\begin{itemize}
\item $x\tr u$ $(1{>}0)$;
\item $\{a,b\}\tr u$ $(2{>}1)$, returning $t_1$;
\item the remaining three of $a,b,c,d$ offered to $u$ $(3{>}2)$, returning $t_2$;
\item $u\tr s_1$ and $u\tr s_2$ for the two remaining rim vertices;
\item $u\tr v$ $(2{>}0)$.
\end{itemize}
For $u\tr s_1$ the supply is $4$ and the blanks are among $x,t_1,t_2$. For $u\tr s_2$ the supply is
$3$, and the blanks are $x$ and at most one of $t_1,t_2$, because $s_2$ has only one rim neighbour.
Finally $b(s_1)=b(s_2)=2$.

\begin{proof}[Proof of Theorem~\ref{thm:main}(b)]
By Lemma~\ref{lem:uv} we may take $u,v$ as above. The contact set $A$ is nonempty and lies in one
of the eleven orbits, so Table~\ref{tab:bowtie} gives a certificate. Lemma~\ref{lem:count} gives a
realisable pair and $p(W)\le(N-1)/2$.
\end{proof}

\begin{remark}\label{rem:bowtie}
(i) Nothing is assumed about $u$ and $v$ beyond the stated adjacencies: they may have any further
neighbours outside $T$, and $v$ need not be a leaf.

(ii) With $u$ alone (a six-vertex local set), certificates exist for $25$ of the $31$ contact sets.
The exceptions are $A\supseteq\{a,b\}$ or $A\supseteq\{c,d\}$ with $A\cap\{c,d\}$ resp.
$A\cap\{a,b\}\in\{\emptyset,\text{both}\}$, that is $ab$, $xab$, $cd$, $xcd$, $abcd$, $xabcd$.

(iii) Random hosts with a bowtie core and no leaf, of orders $9$, $11$ and $13$, all satisfy the
floor bound ($120$ hosts, \texttt{scripts/random\_hosts.py}). We expect Theorem~\ref{thm:main}(b)
to hold without the hypothesis on $u,v$, but the local certificates need it.

(iv) Without the reduction to orbits, \texttt{scripts/certgen.py} finds a certificate for each of the
three labelled bowties and each of the 31 contact sets ($93$ certificates). It also finds a local
certificate for each of the $38$ labelled rims of maximum degree at most two that are not perfect
matchings. These are the certificates replayed in Lean.
\end{remark}

\section{The corrected Corollary 106}\label{sec:cor106}

\begin{corollary}[Corollary~106 of~\cite{MR}, corrected]\label{cor:106}
Let $G$ be connected, $S$ a nonempty set of vertices, and $U$ the set of unmarked even vertices.
Suppose $G[U]$ has at most one vertex of degree at least four, and that $|U|\le5$ or $G[U]$ has
maximum degree at most three. Then $q(G,S)\le\lfloor(|V(G)|-|S|)/2\rfloor$. In particular, in the
two-even-support failure of Proposition~105 with exactly seven even sites, at least two unmarked
even sites have degree four among the unmarked even sites.
\end{corollary}

\begin{proof}
The proof in~\cite{MR} uses its first paragraph, the leaf floor bound, only for graphs $L$ whose
even set is $U$: $L=L_A(G)$ for the nonempty set $A$ of even marks, and $L=G$ with two leaves at an
odd mark. In each case every leaf support is odd. Replace that paragraph as follows.
\begin{itemize}
\item If $|L|$ is even: ceiling and floor agree. Use~\cite[Corollary~14]{BS} when $L[U]$ has maximum
degree at most three, and~\cite[Theorem~5]{FHZ} when it has exactly one vertex of degree at least
four.
\item If $|L|$ is odd: use~\cite[Theorem~13]{BS} when $L[U]$ is subcubic (the leaf excludes SET
graphs). Otherwise $|U|=5$ by hypothesis, and Corollary~\ref{cor:bp5} applies.
\end{itemize}
Replace the citation of~\cite{Shaik} in the case $A=\emptyset$ by~\cite[Corollary~14]{BS}
or~\cite[Theorem~5]{FHZ}, which give the ordinary ceiling for $G$. The rest of the proof is unchanged.
\end{proof}

\begin{corollary}
Theorem~144 of~\cite{MR} holds without~\cite{Shaik}. So does the remark after Theorem~114 that, at
$d_J(z)=6$, a failure of either deletion test requires two universal vertices in the five-vertex
graph $K$.
\end{corollary}

\begin{proof}
By Section~\ref{sec:needs}, every use satisfies the hypothesis of Corollary~\ref{cor:106}, or is
covered by~\cite[Theorem~13]{BS}. In the remark, both test graphs have odd order, leaves at the odd
support $r$, and the five even vertices $N(z)\setminus\{r\}$ (Corollary~\ref{cor:bp5}).
\end{proof}

\emph{What remains conditional.} The leaf floor bound with exactly one exceptional vertex, $|L|$ odd
and $|U|\ge7$ is not proved here, and Proposition~\ref{prop:gadget} shows that the doubling
gadgets of the notes cannot supply it in general. No result of~\cite{MR} uses this case.

\section{Formalization and computations}\label{sec:lean}

\subsection*{Lean}
Two files in \texttt{lean/GallaiLean/} are imported by \texttt{GallaiLean.lean}. They contain no
\texttt{sorry}, and \texttt{\#print axioms} shows only \texttt{propext}, \texttt{Classical.choice} and
\texttt{Quot.sound}. Everything is in the namespace \texttt{Gallai.FiveEven}.

\texttt{FiveEven.lean} (generic):
\begin{itemize}
\item \texttt{okS} and \texttt{soundS}: a version of the insertion-schedule checker \texttt{Prog.ok}
of \texttt{Schedule.lean} whose leaves need not share one end vector, with its soundness theorem.
Each offer is one application of Fan's grouped lemma (\texttt{Ach.grouped}), and the sum of ends is
preserved (\texttt{sum\_mv}).
\item \texttt{seed}: the Lov\'asz seed of a certificate, from \texttt{lovasz\_ends} on
$W-H$ (\texttt{delE}, \texttt{degree\_delE}).
\item \texttt{floor\_of\_cert}: a checked program on a model of $W$, together with the exact model
adjacency and the parities, gives a partition $D$ of $W$ with $2|D|+1=|V|$.
\end{itemize}

\texttt{FiveEvenCerts.lean} (generated by \texttt{scripts/gen\_lean.py} from the certificate tables):
\begin{itemize}
\item the labelled models \texttt{madj5} and \texttt{madj7}, and the tables \texttt{localTable}
($38$ entries) and \texttt{bowTable} ($93$ entries);
\item \texttt{local\_ok} and \texttt{bow\_ok}: every certificate checks, for all $2^6$ rims, resp.
$2^{11}$ rims and contact sets, that satisfy the hypotheses (kernel evaluation,
\texttt{decide +kernel});
\item \texttt{five\_even\_local}, which is Theorem~\ref{thm:main}(a). Let $\iota:\mathrm{Fin}\,5\to V$
be injective onto the even vertices. Suppose $\iota0$ is adjacent to the others, no other $\iota i$ is
adjacent to all the others, and the rim is not a perfect matching. Then $W$ has a partition $D$ with
$2|D|+1=|V|$;
\item \texttt{five\_even\_bowtie}, which is Theorem~\ref{thm:main}(b), given the adjacent odd
vertices $u=\iota5$ and $v=\iota6$ as hypotheses.
\end{itemize}

Both statements quantify over labelled rims, so the classification of rims up to isomorphism is
not needed. Not formalized: Lemma~\ref{lem:uv} (shortest paths), and the reductions of
Corollaries~\ref{cor:bp5} and~\ref{cor:106} to Theorem~\ref{thm:main}. The leaf at an even support in
Corollary~\ref{cor:bp5} is \texttt{rooted\_four}.

\subsection*{Scripts} (\texttt{manuscripts/bs-prescribed/scripts})
\begin{itemize}
\item \texttt{certgen.py} generates the certificates, and \texttt{verify\_certs.py} checks them
independently, in the semantics of the Lean checker: $38$ local certificates with $46$ leaves, and
$93$ bowtie certificates with $129$ leaves.
\item \texttt{bowtie\_table.py} prints Table~\ref{tab:bowtie} with all supply and blank counts.
\item \texttt{sevensite\_cert.py} re-derives audit A6's seven-site bank: every one of the seven cores
and all $31$ contact sets are certified.
\item \texttt{exact\_prescribed.py} and \texttt{odd\_fan\_game.py} produce Table~\ref{tab:census}.
\item \texttt{random\_hosts.py} is an end-to-end consistency check on random hosts, using the circuit
solver of \texttt{tools/pathnum2.py}. All $100$ hosts with a leaf, of orders $9$ and $11$, and all
$120$ bowtie hosts without a leaf satisfy the floor bound.
\item \texttt{sample\_prescribed.py} is the random sample of the prescribed-end lemma at $N=9$
($150$ graphs, $596$ pairs). Its only exceptions are the $6$ pairs with $e=1$ and $3$ pairs with $e=3$.
The Fan game with at most three held edges loses $10$ further pairs that are not exceptions.
\end{itemize}

\begin{thebibliography}{99}
\bibitem{BS} F.~Botler and M.~Sambinelli, Towards Gallai's path decomposition conjecture, J. Graph
Theory 97 (2021) 161--184 (arXiv:1911.04546).
\bibitem{ES} J.~Williams, Endpoint sets of path decompositions at the Gallai budget, working draft,
\texttt{manuscripts/endpoint-structure}.
\bibitem{Fan} G.~Fan, Path decompositions and Gallai's conjecture, J. Combin. Theory Ser. B 93 (2005)
117--125.
\bibitem{FHZ} G.~Fan, J.~Hou and C.~Zhou, Gallai's conjecture on path decompositions, J. Oper. Res.
Soc. China 11 (2023) 439--449.
\bibitem{LF} J.~Williams, Lov\'asz's path theorem from Fan's lemma, \texttt{manuscripts/lovasz-fan}.
\bibitem{L68} L.~Lov\'asz, On covering of graphs, in: Theory of Graphs (Tihany, 1966), Academic
Press, 1968, 231--236.
\bibitem{MR} The marked-reduction manuscript, \texttt{gallai/01-original-manuscripts/marked-reduction.pdf}.
\bibitem{Shaik} I.~A.~Shaik, Gallai's conjecture with two exceptional even vertices, Hexagon preprint,
version~1, 16 October 2026 (unrefereed).
\end{thebibliography}

\end{document}
